% Chapter 6: Conclusion — summary of contributions, limitations, future work
\selectlanguage{greek}

\chapter{Συμπέρασμα}
\label{ch:conclusion}

\section{Σύνοψη Συνεισφορών}
\label{sec:summary}

Η εργασία εισήγαγε το \en{CascadeLP}, μια τετραεπίπεδη αρχιτεκτονική πρόβλεψης
ακμών που δρομολογεί κάθε ζεύγος μέσω προοδευτικά ακριβότερων οικογενειών
χαρακτηριστικών.
Η αρχιτεκτονική αξιολογήθηκε σε δύο σύνολα δεδομένων με
διαφορετική κατασκευή.
Η διπλή αυτή αξιολόγηση αποδείχθηκε ουσιώδης: το πρώτο
σύνολο παρείχε σύγκριση με τον διαγωνισμό \en{DSAA 2023}, ενώ το δεύτερο επέτρεψε
να ελεγχθεί η συμπεριφορά του πλαισίου χωρίς τη συγκεκριμένη στρέβλωση στη
δειγματοληψία που εντοπίστηκε στο πρώτο.

Στο \en{DSAA 2023}, το \en{CascadeLP} πέτυχε \en{Macro F1} =
\CascadeValFone{} υπό το αρχικό πρωτόκολλο επικύρωσης και
\KaggleBaselinePublic{} στο δημόσιο \en{leaderboard}.
Οι τιμές αυτές είναι
έγκυρες ως βαθμολογίες στο συγκεκριμένο πεδίο αξιολόγησης, αλλά δεν αποτελούν
από μόνες τους απόδειξη σχεσιακής πρόβλεψης.
Ο έλεγχος επικάλυψης βρήκε μόνο
\LeakageExact{} ακριβή και \LeakageReversed{} αντεστραμμένα ζεύγη σε
\TestPairs{} ζεύγη ελέγχου, όμως ο ευρύτερος έλεγχος ακεραιότητας αποκάλυψε ότι
\HubCount{} από τις
\HubUniqueSourceCount{} μοναδικές τιμές του πρώτου άκρου στο \en{DSAA 2023}
αντιστοιχούν στο \HubRowPct\% των γραμμών εκπαίδευσης και είναι
\HubPurityPct\% αρνητικές, ενώ κάθε άλλη γραμμή είναι θετική.
Ένας απλός
πίνακας αναζήτησης σε αυτές τις \HubCount{} τιμές πετυχαίνει \en{Macro F1} =
\HubLookupFone{}, και οι προβλέψεις του \en{CascadeLP} συμφωνούν μαζί του στο
\HubValAgreementPct\% του συνόλου επικύρωσης.
Το \en{DSAA 2023} δεν μπορεί συνεπώς να διακρίνει την
εκμάθηση σχέσεων από την αναγνώριση ενός μικρού συνόλου πηγών.
Το ίδιο όριο
ισχύει για τις προηγούμενες σχεδόν τέλειες λύσεις του διαγωνισμού: οι
δημοσιευμένες βαθμολογίες τους δεν αποδεικνύουν ποιο από τα δύο φαινόμενα
έμαθαν τα μοντέλα.

Στο δεύτερο σύνολο πραγματικών υπερσυνδέσμων της Απλής Αγγλικής
\en{Wikipedia}, το ίδιο αρχιτεκτονικό σχήμα επανεκπαιδεύτηκε από την αρχή με
ομαδοποιημένη διαίρεση και απόκρυψη της ακμής-στόχου.
Οι δομικοί ευρετικοί
δείκτες πέτυχαν \en{Macro F1} = \WfStructuralFone{}, ενώ το \en{CascadeLP}
πέτυχε \WfCascadeFone{}, υψηλότερα από κάθε μεμονωμένη οικογένεια
χαρακτηριστικών.
Το \WfTierOnePct\% των ζευγών επιλύθηκε στο Επίπεδο~1, το
\WfTierTwoPct\% στο Επίπεδο~2 και το \WfTierThreePct\% έφτασε στο Επίπεδο~3.
Αυτό αποτελεί την ισχυρότερη ένδειξη της εργασίας ότι η αρχιτεκτονική συνδυάζει
χρήσιμα τις οικογένειες χαρακτηριστικών και δρομολογεί τη μεγάλη πλειονότητα
των ζευγών στο φθηνότερο επίπεδο όταν υπάρχει πραγματικό δομικό σήμα.
Στον συμπληρωματικό έλεγχο με αραιό γράφημα, η επίδοση στα δύσκολα αρνητικά
πέφτει σε \WsRandomHardCascadeFone{} όταν η εκπαίδευση χρησιμοποιεί μόνο
τυχαία αρνητικά. Η μικτή εκπαίδευση την αυξάνει σε
\WsMixedHardCascadeFone{}, υψηλότερα από το αυτόνομο
\en{Embedding} (\WsMixedHardEmbeddingFone{}), με κόστος
\WsRandomCost{} μονάδες \en{Macro F1} στο σύνολο τυχαίων αρνητικών.

Τα τρία ερευνητικά ερωτήματα απαντώνται επομένως με διαφορετικό βαθμό
βεβαιότητας.
Πρώτον, στο ανεξάρτητο σύνολο το \en{CascadeLP} υπερβαίνει την
ισχυρότερη αυτόνομη οικογένεια κατά \WfCascadeDeltaStructural{} μονάδες
\en{Macro F1}, με 95\% διάστημα από επαναδειγματοληψία \en{paired bootstrap}
[\WfCascadeDeltaStructuralCiLow{}, \WfCascadeDeltaStructuralCiHigh{}]. Δεύτερον,
η δρομολόγηση επιλύει το \WfTierOnePct\% των ζευγών στο δομικό επίπεδο και μόνο
το \WfTierThreePct\% φτάνει στο τελευταίο, στοιχείο που τεκμηριώνει την επιλεκτική
κλιμάκωση και τη δυνατότητα μελλοντικής εξαγωγής χαρακτηριστικών κατά απαίτηση,
όχι ήδη μετρημένη ισόποση εξοικονόμηση.
Τρίτον, ο έλεγχος του \en{DSAA 2023}
δείχνει ότι οι σχεδόν τέλειες βαθμολογίες του δεν μπορούν να απομονώσουν τη
σχεσιακή μάθηση από τη στρέβλωση στην επιλογή της πηγής των αρνητικών παραδειγμάτων.

Το μοντέλο αναφοράς \en{TF-IDF + SVM} διατηρήθηκε ως κλασική σημασιολογική
σύγκριση.
Η εκπαίδευση \en{SVC} με πυρήνα ακτινικής βάσης κλιμακώνεται
τετραγωνικά έως κυβικά ως προς το πλήθος δειγμάτων, γι' αυτό στο \en{DSAA 2023}
χρησιμοποιήθηκε στρωματοποιημένο υποσύνολο \SvmSampleSize{} ζευγών.
Η σύγκριση
τοποθετεί το \en{CascadeLP} όχι μόνο απέναντι στις επιμέρους οικογένειες
χαρακτηριστικών του, αλλά και απέναντι σε μια καθιερωμένη μέθοδο με δυσμενέστερο
κόστος εκπαίδευσης.

\section{Περιορισμοί}
\label{sec:limitations-cost}

Η εργασία αξιολογεί το \en{CascadeLP} σε δύο σύνολα υπερσυνδέσμων της
\en{Wikipedia}, όχι όμως σε διαφορετικό τύπο \en{TAG}.
Η επανεκπαίδευση στο
δεύτερο σύνολο ελέγχει την αρχιτεκτονική υπό διαφορετική κατασκευή δεδομένων,
όχι μεταφορά μάθησης χωρίς επανεκπαίδευση ούτε γενίκευση σε άλλο πεδίο εφαρμογής.
Επιπλέον, το δεύτερο σύνολο προέρχεται από διάσχιση \en{BFS} με έναν σπόρο.
Το αρχικό του πρωτόκολλο χρησιμοποιεί ομοιόμορφα τυχαία αρνητικά· ο
συμπληρωματικός έλεγχος με αρνητικά δύο βημάτων δείχνει ουσιώδη πτώση της
επίδοσης και μερική ανάκτησή της μετά από μικτή εκπαίδευση.
Το εύρημα καλύπτει συγκεκριμένη κατασκευή δύσκολων ζευγών, όχι κάθε τρόπο
με τον οποίο τα αρνητικά παραδείγματα μπορούν να μεταβάλουν τη δυσκολία.
Η αναφορά βασίζεται σε έναν σπόρο διάσχισης, μία διαίρεση ανά πρωτόκολλο και μία
εκπαίδευση ανά μοντέλο και εκδοχή αρνητικών.
Συνεπώς, δεν ποσοτικοποιεί τη μεταβλητότητα μεταξύ
διαφορετικών σπόρων, διαιρέσεων και επανεκπαιδεύσεων.
Η επαναδειγματοληψία \en{paired bootstrap} των προβλέψεων επικύρωσης και του
δύσκολου συνόλου ελέγχου παρέχει διαστήματα αβεβαιότητας για τις συγκεκριμένες
σταθερές εκτελέσεις, αλλά δεν υποκαθιστά αυτή την πολλαπλή επανεκτέλεση.
Επιπλέον, στο υποσύνολο μηδενικού \en{CN} του δεύτερου συνόλου το \en{Macro F1}
πέφτει σε \WfZeroCnFone{} και η ανάκληση της θετικής κλάσης σε
\WfZeroCnPositiveRecall{}.
Η συνολική υπεροχή του \en{CascadeLP} τεκμηριώνει τον
συνδυασμό των επιπέδων όταν υπάρχει δομικό σήμα, αλλά δεν τεκμηριώνει αξιόπιστη
πρόβλεψη θετικών ακμών όταν αυτό απουσιάζει.
Το λειτουργικό υποσύνολο
\en{cold-start} περιέχει μόνο \WfColdStartCount{} ζεύγη και δεν επαρκεί για
ξεχωριστό ισχυρό συμπέρασμα.

Η μέτρηση αποτελεσματικότητας της παρούσας υλοποίησης αφορά τη δρομολόγηση και
ταξινόμηση πάνω σε ήδη υπολογισμένα χαρακτηριστικά.
Οι ενσωματώσεις
\en{Sentence-Transformer} υπολογίστηκαν πριν από τη δρομολόγηση για όλους τους
διαθέσιμους κόμβους, επομένως το ποσοστό κλήσεων του Επιπέδου~3 δεν ισούται με
ίση ποσοστιαία μείωση του συνολικού κόστους κωδικοποίησης.
Τέτοια εξοικονόμηση
είναι δυνατή σε υλοποίηση με παραγωγή χαρακτηριστικών κατά απαίτηση και κρυφή
μνήμη ανά κόμβο, αλλά δεν μετρήθηκε ενδογενώς στα παρόντα πειράματα.

Το Επίπεδο~3 (\texttt{\en{EmbeddingClassifier}}) παραμένει αναξιόπιστο στο μικρό
υπόλειμμα ζευγών που το φτάνει (\TierThreeAcc{} ακρίβεια σε \TierThreeCount{} ζεύγη).
Τα ζεύγη αυτού του υπολείμματος που δεν διαχωρίζονται από τα απλά κριτήρια ---
ιδίως εκείνα με πλήρες κείμενο και στις δύο πλευρές,
όπου η ακρίβεια είναι ακόμη χαμηλότερη (\HardResCompleteAcc{}) από το υποσύνολο με ελλιπές
κείμενο --- δεν επιλύονται επαρκώς από την ομοιότητα συνημίτονου ενσωματώσεων
\en{sentence-transformer}.
Το αποτέλεσμα δεν αποδεικνύει ποια πρόσθετη μορφή
πληροφορίας απαιτείται για την επίλυσή τους.
Επιπλέον, οι ετικέτες
«τετριμμένου ζεύγους» του χαρακτηρισμού διαχωρισιμότητας προκύπτουν από ένα όριο
εκατοστημορίου ($\alpha = 0{,}01$) στην κατανομή της αρνητικής κλάσης --- μια
σχεδιαστική επιλογή, όχι μια αντικειμενική διαμέριση του συνόλου δεδομένων.
Το πρωτόκολλο μεταφοράς μάθησης εκτός συνόλου δεδομένων
(\S\ref{sec:generalization-protocol})---εκπαίδευση στο ένα σύνολο, αξιολόγηση
χωρίς περαιτέρω εκπαίδευση σε άλλο---δεν εκτελέστηκε.

Η εργασία αντιμετωπίζει σκόπιμα την πρόβλεψη υπερσυνδέσμων ως πρόβλεψη
\textit{μη κατευθυνόμενης συνδεσιμότητας} μετά από συμμετροποίηση, όχι ως
πρόβλεψη της πραγματικής, κατευθυνόμενης σχέσης παραπομπής
(\S\ref{sec:problem-definition})· η επέκταση σε κατευθυνόμενη πρόβλεψη θα
απαιτούσε ανακατασκευή γραφήματος, χαρακτηριστικών και δειγματοληψίας
αρνητικών και παραμένει εκτός εύρους της παρούσας εργασίας.
Στο ίδιο πνεύμα, ο
έλεγχος αντιστροφής άκρων της \S\ref{sec:endpoint-swap} δείχνει ότι ο
\texttt{\en{PosClassifier}} αλλάζει πρόβλεψη στο \SwapLabelChangePct\% των
ζευγών επικύρωσης όταν αντιστρέφεται η σειρά $u, v$ (\SwapNegClassChangePct\%
στα αρχικά αρνητικά έναντι \SwapPosClassChangePct\% στα αρχικά θετικά), λόγω
της μη-συμμετρικής συνένωσης χαρακτηριστικών \en{POS}
(\S\ref{subsec:pos-features}).
Το εύρημα δεν αμφισβητεί τη συμμετρική
διατύπωση του προβλήματος, αλλά αναδεικνύει περιορισμό της συγκεκριμένης
υλοποίησης του Επιπέδου~2, ο οποίος δεν διορθώθηκε στην παρούσα εργασία.
Επιπλέον, ο επαναληπτικός έλεγχος της κατασκευής γραφήματος
(\S\ref{subsec:data-splits}) δείχνει ότι \AuditValPositivesInGraph{} από τα
\AuditValN{} θετικά ζεύγη επικύρωσης του αρχικού πρωτοκόλλου \en{DSAA} έχουν
ήδη συνεισφέρει ακμή στο γράφημα που χρησιμοποιείται για τα δομικά
χαρακτηριστικά· οι σχετικές τιμές \en{Macro F1} χαρακτηρίζονται συνεπώς ως
αξιολόγηση ακμών σε \en{transductive} διαίρεση με κοινούς κόμβους και με
διαρροή των ακμών επικύρωσης στο γράφημα, όχι ως ανεξάρτητη επικύρωση.
Τέλος, η σύγκριση με το \en{SVM} (\S\ref{sec:svm-results})
συμπληρώθηκε με \MatchedNSeeds{} επαναλήψεις ίσου όγκου εκπαίδευσης
(\MatchedSampleSize{} ζεύγη ανά μοντέλο). Το \en{SVM} πέτυχε
\en{Macro F1} = \MatchedTfidfSvmFone{}, έναντι \MatchedCascadeFone{} του
\en{CascadeLP} και \MatchedPosFone{} του αυτόνομου \en{POS}. Η σύγκριση είναι
δίκαιη ως προς τον αριθμό εποπτευόμενων ζευγών και τον προϋπολογισμό
επισημασμένων ζευγών για την κατασκευή του γραφήματος,
αλλά παραμένει εντός του \en{DSAA 2023} και επομένως υπόκειται στη στρέβλωση
στη δειγματοληψία αρνητικών παραδειγμάτων.
Η επανεκπαίδευση από την
αρχή σε ανεξάρτητο σύνολο δεδομένων \en{Wikipedia} (\S\ref{sec:wiki-fresh})
δείχνει ότι η αρχιτεκτονική συμπεριφέρεται όπως σχεδιάστηκε όταν απουσιάζουν τα
ζητήματα ακεραιότητας του \en{DSAA 2023}, αλλά δεν ελέγχει αν αναπαραστάσεις
εκπαιδευμένες στο ένα σύνολο δεδομένων μεταφέρονται σε άλλο χωρίς
επανεκπαίδευση.

\section{Μελλοντική Έρευνα}
\label{sec:future-work}

Η σαφέστερη κατεύθυνση που προκύπτει άμεσα από αυτή την εργασία είναι η
εφαρμογή μιας \textit{inductive} μεθόδου αναπαράστασης κόμβων, όπως το \en{GraphSAGE}, ή γενικότερα με ένα \en{Graph Neural Network (GNN)} που υπολογίζει την αναπαράσταση ενός κόμβου από τα χαρακτηριστικά και τη γειτονιά του.
Ένα τέτοιο \en{GNN} θα μπορούσε, μέσω \en{message passing} πάνω σε χαρακτηριστικά κόμβων (π.χ. \en{TF-IDF} ή \en{sentence-transformer} ενσωματώσεις κειμένου άρθρου), να παραγάγει αναπαράσταση για κόμβο που δεν συμμετείχε σε καμία ακμή εκπαίδευσης και ενδεχομένως να αντικαταστήσει ένα ή περισσότερα επίπεδα της αλυσίδας.

Δεύτερη κατεύθυνση είναι η τυπική μεταφορά μάθησης μεταξύ συνόλων δεδομένων
\en{TAG}---εκπαίδευση στο ένα, αξιολόγηση χωρίς περαιτέρω εκπαίδευση στο
άλλο---που παραμένει ανοιχτή παρά την επανεκπαίδευση από την αρχή σε
ανεξάρτητο σύνολο δεδομένων \en{Wikipedia} της \S\ref{sec:wiki-fresh}, καθώς
και η επέκταση αυτής της επικύρωσης σε συνόλα δεδομένων εκτός \en{Wikipedia}
(π.χ.\ αναφορές επιστημονικών άρθρων, ή γραφήματα κοινωνικών δικτύων με
κειμενικά προφίλ χρηστών), ώστε να ελεγχθεί αν το πρωτόκολλο χαρακτηρισμού
δυσκολίας και η κατανομή κλήσεων ανά επίπεδο του \en{CascadeLP} γενικεύονται
πέρα από γραφήματα υπερσυνδέσμων.
Τρίτη κατεύθυνση είναι η επανάληψη της ανεξάρτητης αξιολόγησης με πολλαπλούς
σπόρους \en{BFS} και διαιρέσεις, καθώς και με δυσκολότερα αρνητικά παραδείγματα
ταιριασμένα ως προς βαθμό, θέμα ή κειμενική ομοιότητα.
Έτσι μπορεί να διαχωριστεί
η επίδραση της αρχιτεκτονικής από τη μεταβλητότητα της κατασκευής του συνόλου.
Τέταρτη κατεύθυνση είναι η επέκταση σε χρονική πρόβλεψη ακμών, αξιοποιώντας τη χρονική
σήμανση των άρθρων \en{Wikipedia} για την πρόβλεψη \textit{πότε} θα εμφανιστεί μια
ακμή και όχι μόνο αν θα εμφανιστεί.

Πέμπτη κατεύθυνση είναι η αντιμετώπιση του περιορισμού σειράς άκρων της
\S\ref{sec:endpoint-swap}: συμμετροποίηση του διανύσματος \en{POS} (π.χ.
άθροισμα ή απόλυτη διαφορά αντί για συνένωση με σειρά), ή εναλλακτικά επέκταση
του προβλήματος σε γνήσια κατευθυνόμενη πρόβλεψη παραπομπής αν το ερώτημα
έρευνας το δικαιολογεί, κάτι που θα απαιτούσε αναδιατύπωση του γραφήματος και
των δομικών χαρακτηριστικών (\S\ref{sec:problem-definition}).

Τέλος, η διαχείριση κόμβων με ελλιπές κείμενο άρθρου --- υπεύθυνων για το \HardResMissingPct\% του
\en{hard residual} του Επιπέδου~3, χωρίς να αποτελούν την κύρια αιτία της
αποτυχίας του (\S\ref{sec:hard-residual}) --- αξίζει ξεχωριστή διερεύνηση: πιθανές
προσεγγίσεις περιλαμβάνουν ανάκτηση εναλλακτικής πηγής κειμένου (π.χ.\ σύνοψη από
συνδεδεμένα άρθρα) ή ρητή σήμανση της απουσίας κειμένου ως χαρακτηριστικό, ώστε η
αλυσίδα να αναγνωρίζει αυτά τα ζεύγη ως δομικά διαφορετική κατηγορία αντί να τα
αντιμετωπίζει ως συνηθισμένη σημασιολογική αποτυχία.
